\documentstyle[% 


righttag% 


,11pt% 


,amssymb% 


,russian%               remove in Eng. 


,izvru%                 Set izven% in Eng. 


]{amsart} 


 


%  


 


\newcommand{\la}{\langle} 


\newcommand{\ra}{\rangle} 


\newcommand{\supp}{\operatorname{supp}} 


\newcommand{\re}{\operatorname{Re}} 


\newcommand{\sgn}{\operatorname{sgn}} 


\newcommand{\tts}[1]{} 


\renewcommand{\baselinestretch}{0.96}


 


\theoremstyle{plain} 


\theorembodyfont{\it} 


\newtheorem{thms}{\hskip\parindent Теорема}[section] 


\newtheorem{lems}{\hskip\parindent Лемма}[section] 


\theoremstyle{definition} 


\theorembodyfont{\rm} 


\newtheorem{notes}{\hskip\parindent Замечание}[section] 


\newtheorem{cors}{\hskip\parindent Следствие}[section] 


\numberwithin{equation}{section} 


\setcounter{section}{5}  


\setcounter{equation}{14}  


\setcounter{notes}{2} 


\setcounter{thms}{1} 


\setcounter{lems}{3} 


 


%\hoffset=-15mm%                  For Eng. 


  


\setcounter{page}{50} 


\god{1999} 


\nomer{11~(450)} %                        Remove in Eng. 


%\nomer{Vol.\,43, No.\,11}%               Set in Eng. 


%\str{\pageref{begin}--\pageref{end}}%  Set in Eng. 


\udk{517.9} 


 


\author{\it В.С.\,Мокейчев, \fbox{А.В.\,Мокейчев}}   


% NB:   ^^ remove in Eng.   


\title{Новый подход к теории линейных задач\\   


для систем дифференциальных уравнений\\   


в частных производных, III} 


 


\date{} 


% \pagestyle{myheadings}%         Set in Eng. 


\pagestyle{plain} %              Remove in Eng. 


\begin{document} 


% \thispagestyle{plain}%          Set in Eng. 


\maketitle 


\label{begin} 


 


В данной статье сохраняются обозначения, понятия, соглашения и 


нумерация, использованные в первых двух частях \cite{1}, \cite{2}. 


Рассмотрим второе применение ранее полученных абстрактных результатов. 


\medskip 


 


2. $(b{-}a)$-{\it периодическая задача}. В этом пункте 


$H=L^2_m((a,b))$, $\varphi$ совпадает с (1.4), т.\,е, 


$$ 


\varphi=\{e(j)\exp(2\pi i k\bullet\tfrac{x-a}{b-a}, 


\ j=1,\dotsc,m,\ k\in\Bbb Z^n\}. 


$$ 


Ради простоты считаем, что при каждом $\alpha\in\Phi$ оператор 


$C_\alpha(x)$ является оператором умножения на матрицу $C_\alpha(x)$, 


являющуюся измеримой и квадратной. Размерности матриц $(m,m)$, 


$m<+\infty$. 


 


\begin{thms} 


Пусть $C_\alpha(x)\in C^{\{\alpha\}}_{m,m}(b-a)$ 


$\forall\,\alpha\in\Phi$, 


$q(z)=\sum\limits_{\alpha\in\Phi}a_\alpha(iz)^\alpha$ --- скалярный 


многочлен, удовлетворяющий условиям 


\begin{equation} 


|q(k)|^2\ge g_1\sum_{\alpha\in\Phi}\sum_{\beta\le\alpha} 


|k^\beta|^2,\quad 


\text{как только} 


\quad k\in\Bbb Z^n, 


\end{equation} 


при всех $\gamma<\alpha$, каждом $\alpha\in\Phi$, всех $|k|\ge 


N(\varepsilon_1)$ и некотором $\varepsilon_1>0$ 


\begin{equation} 


(|k^\gamma|/|q(k)|)<\varepsilon_1, 


\end{equation} 


причем постоянные $g_1>0$, $N(\varepsilon_1)$ и $\varepsilon_1$ не 


зависят от $k$. Если $\varepsilon_1$ достаточно мало, то следующие 


утверждения равносильны$:$ 


\begin{itemize} 


\item[1.] 


существуют такие числа $t\in\Bbb R$, $d_1(t)$, $d_2>0$, 


$\alpha_0\in\Bbb C$, что при каждом $x\in[a,b]$ и всех $z\in\Bbb C^m$, 


$k\in\Bbb Z^n$, выполняются оценки 


\begin{equation} 


\bigg|\bigg(\alpha_0t E+\sum_{\alpha\in\Phi}C_\alpha(x) 


\bigg(\frac{2\pi i k}{b-a}\bigg)^\alpha\bigg)z\bigg|^2\ge 


|z|^2(d_1(t)+d_2|q(k)|^2), 


\end{equation} 


в которых $d_1(t)\to+\infty$, если $|t|\to+\infty;$ 


\item[2.] 


существуют такие числа $t\in\Bbb R$, $d_3(t)$, $d_4>0$, $d_5>0$, 


$\alpha_0\in\Bbb C$, что при всех $h(x)\in C^{+\infty}_m(b-a)$ 


выполняются оценки 


\begin{multline} 


\|(P(x,\partial/\partial x)+\alpha_0t E)h(x)\|\ge d_3(t) 


\|h(x)\|+d_4\|q(-i\partial/\partial x)h(x)\|\ge \\ 


% 


\ge d_3(t)\|h(x)\|+d_5\sum_{\alpha\in\Phi} 


\sum_{\beta\le\alpha}\|h^{(\beta)}(x)\|, 


\end{multline} 


в которых $d_3(t)\to+\infty$ в случае $|t|\to+\infty$. 


\end{itemize} 


\end{thms} 


 


 


\begin{notes} 


Если $q(-i\xi)$ --- гипоэллиптический полином, то (5.16) имеет место 


при каждом $\varepsilon_1>0$ и всех $|k|\ge N(\varepsilon_1)$. Однако 


(5.16) выполняется при малых $\varepsilon_1$ не только для 


гипоэллиптических полиномов. Доказывая теорему 5.3, убедимся в этом. 


\end{notes} 


 


\begin{notes} 


В частном случае $m=1$ и при более жестких ограничениях на $q(k)$ 


теорема 5.2 доказана в (\cite{3}, с.\,141--152). 


\end{notes} 


 


Большую роль будет играть 


 


\begin{lems} 


Пусть множество $\Omega$ ограничено, элементы матриц $\wt 


C_{\alpha,\beta}(x)$ ограничены в существенном и имеет место ($5.16$). 


Тогда 


\begin{equation} 


S_4=\bigg(\int_\Omega\bigg|\sum_{\alpha\in\Phi} 


\sum_{\beta<\alpha}\wt C_{\alpha,\beta}(x) 


h^{(\beta)}(x)\bigg|^2dx\bigg)^{1/2}\le 


\varepsilon_1B\|q(-i\partial/\partial x)h(x)\|+c\|h(x)\|, 


\end{equation} 


причем постоянные $B$, $c$ не зависят от $h(x)$, и $B$ не зависит от 


$\varepsilon_1$. 


\end{lems} 


 


\begin{pf} 


Так как $|\wt C_{\alpha,\beta}(x)|\le\wt B$ почти всюду, то 


$$ 


S_4\le\wt B\sum_{\alpha\in\Phi}\sum_{\beta<\alpha} 


\|h^{(\beta)}(x)\|. 


$$ 


Напомним, что $\|\cdot\|$ --- символ нормы в $L^2_m((a,b))$. Из 


последних оценок и из (5.16) следует неравенство 


\begin{equation} 


S_4\le B_1\varepsilon_1\bigg(\sideset{}{'}\sum|h_k|^2 


|q(k)|^2\bigg)^{1/2}+B\sum_{\alpha\in\Phi} 


\sum_{\beta<\alpha}\bigg(\sideset{}{''}\sum|h_k|^2 


|k^\beta|^2\bigg)^{1/2}. 


\end{equation} 


Здесь штрих указывает на то, что суммирование производится по всем 


$k\in\Bbb Z^n\cap(|k|\ge N(\varepsilon_1))$, два штриха --- по всем 


$k\in\Bbb Z^n\cap(|k|<N(\varepsilon_1))$, $N(\varepsilon_1)$ --- 


постоянная, используемая в (5.16), и число $B_1$ равно произведению 


$\wt B$ на число элементов в множестве 


$\mathop{\cup}\limits_{\alpha\in\Phi}\{M^n\cap[0,\alpha)\}$. Из (5.20) 


легко следует (5.19) 


\end{pf} 


 


\begin{lems} 


Пусть $\Omega_j=\{x:|x-\tau(j)|<\rho_j\}$, 


$\Omega=\Omega_1\cup\dotsb\cup\Omega_M$, и $[a,b]\subset\Omega$. Тогда 


существуют такие вещественные скалярные функции $v_j(x)\in 


C^{+\infty}_0(\Omega)$, $j=1,\dotsc,M$, что 


\begin{equation} 


0\le v_j(x)\le 1\ \; \forall\,j;\quad 


(\supp v_j(x))\subset\ov\Omega\ \; \forall\,j;\quad 


(v_1(x))^2+\dotsb+(v_M(x))^2=1\ \; \forall\,x\in[a,b]. 


\end{equation} 


\end{lems} 


 


\begin{pf} 


Напомним, что $(\supp v_j(x))$ --- носитель функции $v_j(x)$ и 


$\ov\Omega_j$ --- замыкание множества $\Omega_j$. Лемма доказывается 


аналогично тому, как доказывается существование классического разбиения 


единицы. Положим 


$\wt\omega_j(x)=\exp((|x-\tau(j)|^2-\rho^2_j)^{-1})$ при $x\in\Omega_j$, 


$\wt\omega_j(x)=0$


если $x\notin\Omega_j$. Через $\wt\omega_0(x)\in C^{+\infty}_0(\Omega)$ 


обозначим функцию, удовлетворяющую условию $\wt\omega_0(x)=1$ при 


$x\in[a,b]$. Легко заметить, что функции 


$$ 


v_j(x)=\wt\omega_0(x)\wt\omega_j(x)\big((\wt\omega_1(x))^2+ 


\dotsb+(\wt\omega_M(x))^2\big)^{1/2} 


$$ 


удовлетворяют (5.21). 


\end{pf} 


 


\begin{notes} 


Сформулировать и доказать лемму 5.5 было необходимо по той причине, 


что из включения $v_0(x)\in C^{+\infty}_0(\Omega)$ не следует 


$\sqrt{v_0(x)}\in C^{+\infty}_0(\Omega)$. 


\end{notes} 


 


{\bf Доказательство теоремы 5.2}. Фиксируем достаточно малое число 


$\varepsilon>0$ и числа $N(\varepsilon)$, $\rho(\varepsilon)$ так, 


чтобы 


\begin{equation} 


|C_\alpha(x)-C_\alpha(y)|<\varepsilon\quad 


\text{при}\ \; x\in\Omega_j,\ \; y\in\Omega_j,\ \; 


j=1,\dotsc,N(\varepsilon), 


\end{equation} 


где $\Omega_j=\{x:|x-\tau(j)|<\rho(\varepsilon)\}$ и 


$[a,b]\subset\Omega\equiv\Omega_1\cup\dotsb\cup\Omega_{N(\varepsilon)}$


. Пусть $v_j(x)$, $j=1,\dotsc,N(\varepsilon)$, --- функции, построенные 


в лемме 5.5. Так как $C_\alpha(x)$ и $h(x)$ --- 


$(b{-}a)$-периодические, то при некоторой постоянной $B_2$, не 


зависящей от $t$, $h(x)$, 


$$ 


S_5\equiv\bigg(\int_\Omega|(P(x,\partial/\partial x)+ 


\alpha_0t E)h(x)|^2dx\bigg)^{1/2}\le B_2 


\|(P(x,\partial/\partial x)+\alpha_0 t E)h(x)\|. 


$$ 


С другой стороны, имеет место (5.21), поэтому 


\begin{multline*} 


S_5\ge\bigg(\sum_j\int_\Omega|v_j(x)(P(x,\partial/\partial x)+ 


\alpha_0t E)h(x)|^2dx\bigg)^{1/2} \ge \\ 


% 


\ge\bigg(\sum_j\int_\Omega(P(\tau(j),\partial/\partial x)+ 


\alpha_0t E)(v_j(x)h(x))|^2dx\bigg)^{1/2}- \\ 


% 


-\bigg(\sum_j\int_\Omega|(P(x,\partial/\partial x)- 


P(\tau(j),\partial/\partial x))(v_j(x)h(x))|^2dx\bigg)^{1/2}- \\ 


% 


-\bigg(\sum_j\int_\Omega|v_j(x)P(x,\partial/\partial x) 


h(x)-P(x,\partial/\partial x)(v_j(x)h(x))|^2dx\bigg)^{1/2}. 


\end{multline*} 


Слагаемое в правой части обозначим через $F_4$, вычитаемые --- 


соответственно через $F_5$, $F_6$. Оценим $F_5$, $F_6$ сверху, а $F4$ 


--- снизу. В силу обычных неравенств для норм и (5.22) 


\begin{multline*} 


F_5\le\varepsilon\bigg(\sum_j\sum_{\alpha\in\Phi}\int_\Omega 


|(v_j(x)h(x))^{(\alpha)}|^2dx\bigg)^{1/2}\le\varepsilon\bigg( 


\sum_j\sum_{\alpha\in\Phi}\int_\Omega|v_j(x)h^{(\alpha)}(x)|^2 


dx\bigg)^{1/2}+ \\ 


% 


+\varepsilon\bigg(\sum_j\sum_{\alpha\in\Phi}\int_\Omega 


|v_j(x)h^{(\alpha)}(x)-(v_j(x)h(x))^{(\alpha)}|^2dx\bigg)^{1/2}. 


\end{multline*} 


Группы слагаемых в правой части последней оценки обозначим 


соответственно через $F_{5,1}$, $F_{5,2}$. Так как имеет место (5.21), 


то 


\begin{equation} 


F_{5,1}\le\varepsilon\bigg(\sum_{\alpha\in\Phi}\int_\Omega 


|h^{(\alpha)}(x)|^2dx\bigg)^{1/2}\le\varepsilon B_1 


\sum_{\alpha\in\Phi}\|h^{(\alpha)}(x)\|. 


\end{equation} 


Напомним, что $\|\cdot\|$ --- символ нормы в $L^2_m((a,b))$. Применяя 


лемму 5.1, получим 


\begin{equation} 


F_{5,2}\le\varepsilon_1 B\|q(-i\partial/\partial x)h(x)\|+ 


c\|h(x)\|. 


\end{equation} 


В силу (5.15) имеем 


\begin{equation} 


\sum_{\alpha\in\Phi}\|h^{(\alpha)}(x)\|\le g_2 


\|q(-i\partial/\partial x)h(x)\|+c\|h(x)\|. 


\end{equation} 


Из (5.23), (5.24), (5.25) следует 


\begin{equation} 


F_5\le(\varepsilon B_1g_2+\varepsilon_1B)\| 


q(-i\partial/\partial x)h(x)\|+c\|h(x)\|, 


\end{equation} 


причем постоянные $\varepsilon$, $\varepsilon_1$, $B$, $B_1$, $g_2$, 


$c$ не зависят от $t$, $h(x)$; постоянные $B_1$, $g_2$ не зависят от 


$\varepsilon$, $\varepsilon_1$; $B$, $\varepsilon$ не зависят от


$\varepsilon_1$. Чтобы оценить $F_6$, достаточно воспользоваться леммой 


5.4 


\begin{equation} 


F_6\le(\varepsilon_1B)\|q(-i\partial/\partial x)h(x)\|+c\|h(x)\|. 


\end{equation} 


 


Труднее оценить $F_4$. Предварительно оценим 


$$


F_{4,j}=\int_\Omega|(P(\tau(j),\partial/\partial x)+ 


\alpha_0t E)(v_j(x)h(x))|^2dx =


\int_{\Omega_j}|(P(\tau(j),\partial/\partial x)+ 


\alpha_0t E)(v_j(x)h(x))|^2dx. 


$$


Уменьшая, если в этом есть необходимость, $\rho(\varepsilon)$, можно 


добиться того, что $\Omega_j\subset[a+\tau(j),\;b+\tau(j)]$. В этом 


случае при некоторой $\wt v_j(x)\in C^{+\infty}_1(b-a)$ и всех 


$x\in[a+\tau(j),\;b+\tau(j)]$ имеют место равенства $\wt 


v_j(x)=v_j(x)$, 


$$ 


F_{4,j}=\int^{b+\tau(j)}_{a+\tau(j)} 


|(P(\tau(j),\partial/\partial x)+\alpha_0t E) 


(\wt v_j(x)h(x))|^2dx. 


$$ 


Для $(b{-}a)$-периодических функций $f(x)$ выполняется равенство 


\begin{equation} 


\int^b_a|f(x)|^2dx=\int^{b+\tau(j)}_{a+\tau(j)} 


|f(x)|^2dx 


\end{equation} 


(естественно при условии существования интеграла). Поэтому 


$$ 


F_{4,j}=\int^b_a|(P(\tau(j),\partial/\partial x) 


+\alpha_0t E)(\wt v_j(x)h(x))|^2dx= 


\sum_k\bigg|\bigg(P\bigg(\tau(j),\frac{2\pi i k}{b-a}\bigg) 


+\alpha_0t E\bigg)z_k\bigg|^2. 


$$ 


Здесь $z_k$ --- вектор-столбец с координатами $z_{k,r}$, являющимися 


коэффициентами Фурье с номерами $(k,r)$ для функции 


$(\wt v_j(x)h(x))$. Напомним, что $\varphi$ совпадает с (1.4). Из 


последних равенств и из (5.17) следует 


\begin{equation} 


F_{4,j}\ge\sum_k|z_k|^2\bigg(d_1(t)+d_2\bigg|q\bigg( 


\frac{2\pi l}{b-a}\bigg)\bigg|^2\bigg)=d_1(t) 


\|\wt v_j(x)h(x)\|^2+d_2\|q(-i\partial/\partial x) 


(\wt v_j(x)h(x))\|^2. 


\end{equation} 


Учитывая (5.28), получим 


\begin{multline} 


\|\wt v_j(x)h(x)\|^2=\int^{b+\tau(j)}_{a+\tau(j)} 


|\wt v_j(x)h(x)|^2dx=\int^{b+\tau(j)}_{a+\tau(j)} 


|v_j(x)h(x)|^2dx\ge \\ 


% 


\ge\int_\Omega|v_j(x)h(x)|^2dx\ge\|v_j(x)h(x)\|^2. 


\end{multline} 


При этом были учтены включения 


$(\supp v_j(x))\subseteq\ov\Omega_j\subseteq 


[a+\tau(j),\;b+\tau(j)]$. 


 


Оценим снизу величину 


\begin{multline*} 


\|q(-i\partial/\partial x)(\wt v_j(x)h(x))\|^2\ge\frac{1}{2} 


\|\wt v_j(x)q(-i\partial/\partial x)h(x)\|^2- \\ 


% 


-\frac{1}{2}\|q(-i\partial/\partial x)(\wt v_j(x)h(x))- 


\wt v_j(x)q(-i\partial/\partial x)h(x)\|^2. 


\end{multline*} 


Аналогично тому, как было получено (5.30), получим 


$$ 


\|\wt v_j(x)q(-i\partial/\partial x)h(x)\|^2\ge 


\|v_j(x)q(-i\partial/\partial x)h(x)\|^2. 


$$ 


С другой стороны, в силу леммы 5.4 


$$ 


\|q(-i\partial/\partial x)(\wt v_j(x)h(x))-\wt v_j(x) 


q(-i\partial/\partial x)h(x)\|\le \varepsilon_1B 


\|q(-i\partial/\partial x)h(x)\|+c\|h(x)\|. 


$$ 


Таким образом, доказано, что 


$$ 


\|q(-i\partial/\partial x)(\wt v_j(x)h(x))\|^2\ge\frac{1}{2} 


\|v_j(x)q(-i\partial/\partial x)h(x)\|^2- 


\varepsilon_1B\|q(-i\partial/\partial x)h(x)\|^2- 


c\|h(x)\|^2. 


$$ 


Отсюда и из (5.29), (5.30) следует неравенство 


\begin{equation} 


F_{4,j}\ge d_1(t)\|v_j(x)h(x)\|^2+d_2(\tfrac{1}{2} 


-\varepsilon_1B)\|v_j(x)q(-i\partial/\partial x)h(x)\|^2- 


c\|h(x)\|^2. 


\end{equation} 


Так как $c$ не зависит от $t$ и $d_1(t)\to+\infty$ при некоторых 


$|t|\to+\infty$, то можно считать $d_1(t)-c\ge0$. Число $\varepsilon_1$ 


будем считать настолько малым, что $\varepsilon_1B<\frac{1}{2}$. При 


сделанных оговорках правая часть в (5.31) неотрицательна. Из (5.31) и 


из обычных неравенств для норм следует 


\begin{equation} 


F_4\ge\big(d_2(\tfrac{1}{2}-\varepsilon_1B)\big)^{1/2} 


\|q(-i\partial/\partial x)h(x)\|+ 


(d_1(t)-c)^{1/2}\|h(x)\|. 


\end{equation} 


Учитывая (5.26), (5.27), (5.32), получим 


\begin{multline} 


\|(P(x,\partial/\partial x)+\alpha_0t E)h(x)\|\ge\Big(\big(d_2(\tfrac{1}{2}- 


\varepsilon_1B)\big)^{1/2}-\varepsilon B_1g_2- 


2\varepsilon_1B\Big)\|q(-i\partial/\partial x)h(x)\|+ \\


%


+\big((d_1(t)-c)^{1/2}-c\big)\|h(x)\|, 


\end{multline} 


причем все постоянные не зависят от $h(x)$; все постоянные, кроме 


$d_1(t)$, не зависят от $t$; постоянные $B_1$, $g_2$, $d_2$ не зависят 


от $\varepsilon$, $\varepsilon_1$; постоянная $\varepsilon$ не зависит 


от $\varepsilon_1$. Из доказанной оценки (5.33) при достаточно малых 


$\varepsilon$, $\varepsilon_1$ следует первая из оценок (5.18), и с 


учетом (5.25) --- вторая. Импликация $1\,\Rightarrow\,2$ доказана. 


 


Пусть выполняется утверждение 2 и $y\in[a,b]$. Тогда $y\in(\supp 


v_j(x))$ при некотором $j$. Постоянную $\rho(\varepsilon)$, 


используемую в (5.22), фиксируем так, чтобы 


$\Omega_j\subset[a+\tau(j),\;b+\tau(j)]$. Через $\wt v_j(x)$, как и 


выше, обозначим функцию, принадлежащую $C^{+\infty}_1(b-a)$ и 


удовлетворяющую условию $\wt v_j(x)=v_j(x)$ для всех 


$x\in[a+\tau(j),\;b+\tau(j)]$. Если $h(x)\in C^{+\infty}_m(b-a)$, то 


\begin{multline*} 


\|\wt v_j(x)(P(y,\partial/\partial x)+\alpha_0t E)h(x)\|\ge 


\|(P(x,\partial/\partial x)+\alpha_0t E) 


(\wt v_j(x)h(x))\|- \\ 


% 


-\|(P(x,\partial/\partial x)-P(y,\partial/\partial x)) 


(\wt v_j(x)h(x))\|-\|\wt v_j(x) 


P(y,\partial/\partial x)h(x)-P(y,\partial/\partial x) 


(\wt v_j(x)h(x))\|. 


\end{multline*} 


Слагаемое в правой части обозначим через $\wt F_4$, вычитаемые --- 


соответственно $\wt F_5$, $\wt F_6$. Аналогично тому, как были оценены 


$F_5$, $F_6$, используемые в доказательстве импликации 


$1\,\Rightarrow\,2$, получим 


\begin{align*} 


\wt F_5&\le\varepsilon\|q(-i\partial/\partial x) 


(\wt v_j(x)h(x))\|+c\|h(x)\|, \\ 


% 


\wt F_6&\le\varepsilon_1B\|q(-i\partial/\partial x) 


h(x)\|+c\|h(x)\|. 


\end{align*} 


В силу утверждения 1 


$$ 


\wt F_4\ge d_4\|q(-i\partial/\partial x) 


(\wt v_j(x)h(x))\|+d_3(t)\|h(x)\|. 


$$ 


Таким образом, 


\begin{multline*} 


\|\wt v_j(x)(P(y,\partial/\partial x)+\alpha_0t E)h(x)\|\ge \\ 


% 


\ge (d_4-\varepsilon)\|q(-i\partial/\partial x) 


(\wt v_j(x)h(x))\|-\varepsilon_1B 


\|q(-i\partial/\partial x)h(x)\|+(d_3(t)-c)\|h(x)\|. 


\end{multline*} 


В силу леммы 5.4 


$$ 


\|q(-i\partial/\partial x)(\wt v_j(x)h(x))\|\ge 


\|\wt v_j(x)q(-i\partial/\partial x)h(x)\|- 


\varepsilon_1B\|q(-i\partial/\partial x)h(x)\|-c\|h(x)\|. 


$$ 


Из двух последних оценок следует 


\begin{multline*} 


\|\wt v_j(x)(P(y,\partial/\partial x)+\alpha_0t E)h(x)\|\ge \\ 


% 


\ge (d_4-\varepsilon)\|\wt v_j(x) 


q(-i\partial/\partial x)h(x)\|- 


((d_4-\varepsilon)\varepsilon_1B+\varepsilon_1B) 


\|q(-i\partial/\partial x)h(x)\|+(d_3(t)-c)\|h(x)\|. 


\end{multline*} 


Напомним, что мы условились писать несущественные постоянные без 


индексов. Полагая 


$h(x)=z\exp(2\pi ip\bullet\frac{x-a}{b-a})$, где $z\in \Bbb C^m$ не 


зависит от $x$, получим 


\begin{multline*} 


\|\wt v_j(x)\|\bigg|\bigg(P\bigg(y,\frac{2\pi ip}{b-a}\bigg) 


+\alpha_0t E\bigg)z\bigg|\ge (d_4-\varepsilon) 


\|\wt v_j(x)\|\bigg|q\bigg(\frac{2\pi p}{b-a}\bigg)z\bigg|- \\ 


% 


-((d_4-\varepsilon)\varepsilon_1B+\varepsilon_1B)\bigg|q\bigg( 


\frac{2\pi p}{b-a}\bigg)z\bigg|\bigg(\prod^n_{j=1}\bigg( 


\frac{(b_j-a_j)}{2\pi}\bigg)\bigg)^{1/2}+(d_3(t)-c)|z|\bigg( 


\prod^n_{j=1}\bigg(\frac{(b_j-a_j)}{2\pi}\bigg)\bigg)^{1/2}. 


\end{multline*} 


В полученных оценках постоянные не зависят от $h(x)$; все постоянные, 


кроме $d_3(t)$, не зависят от $t$; $d_4$, $d_3(t)$ не зависят от 


$\varepsilon$, $\varepsilon_1$; $B$, $\varepsilon$ не зависят от 


$\varepsilon_1$. Поэтому при достаточно малых $\varepsilon$, 


$\varepsilon_1$ выполняются оценки (5.17). $\quad\square$ 


 


\begin{cors} 


Пусть выполняются предположения теоремы 5.2, оценки (5.17), 


$\varepsilon_1$ достаточно мало, при достаточно малом 


$\varepsilon_2>0$, некотором $N(\varepsilon_2)$ и всех $|k|\ge 


N(\varepsilon_2)$ имеет место $\varepsilon_2|q(\frac{2\pi k}{b-a})|\ge 


1$, оператор $Q$ задается равенствами 


\begin{equation} 


Q v=\sum_k q(k)v_k\exp\bigg(2\pi ik\bullet\frac{x-a}{b-a}\bigg) 


\ \; \forall\,v\in D_\varphi. 


\end{equation} 


Тогда выполняются оценки (2.1), если, кроме того, $C_\alpha(x)\in 


C^{\{\alpha\}}_{m,m}(b-a)$ при всех $\alpha\in\Phi$, то выполняются и 


оценки (2.8). 


\end{cors} 


 


\begin{cors} 


Пусть $C_\alpha(x)\in C^{\{\alpha\}}_{m,m}(b-a)$ при всех 


$\alpha\in\Phi$, выполняются предположения теоремы 5.2, оценки (5.18) и 


равенства (5.34). Тогда при $c_3=g_3\varepsilon_1$ и некоторой 


постоянной $g_3$, не зависящей от $\varepsilon_1$ и $\psi(x)$, имеет 


место (2.13). 


\end{cors} 


 


\begin{cors} 


Предположим, что $C_\alpha(x)\in C^{\{\alpha\}}_{m,m}(b-a)$ при всех 


$\alpha\in\Phi$, выполняются оценки (5.17) и $q(i\xi)$ --- 


гипоэллиптический полином. Тогда выполняются оценки (2.1), (2.8) и при 


любых $c_3>0$ --- оценки (2.13). 


\end{cors} 


 


{\bf Доказательство следствия 5.1.} 


Напомним, что $\varphi$ совпадает с (1.4). Так как 


$$ 


\|P_\lambda\psi\|\ge\|(P(x,\partial/\partial x)+ 


\alpha_0t E)\psi\|-|t\alpha_0+\lambda|\,\|\psi\|, 


$$ 


то в силу теоремы 5.2 


\begin{multline*} 


\|P_\lambda\psi\|\ge d_4\|q(-i\partial/\partial x)\psi\|+ 


d_3(t)\|\psi\|-|t\alpha_0+\lambda|\,\|\psi\| = \\ 


% 


=d_4\|q(-i\partial/\partial x)\psi\|+(d_3(t)- 


|t\alpha_0+\lambda|)\|\psi\|. 


\end{multline*} 


С другой стороны, 


$\|\psi\|\le\varepsilon_2B\|q(-i\partial/\partial x)\psi\|+ 


c\|\psi(N,x)\|$. 


Поэтому 


$$ 


\|P_\lambda\psi\|\ge(d_4-\varepsilon_2B|d_3(t)- 


|t\alpha_0+\lambda|\,|)\|q(-i\partial/\partial x)\psi\|- 


c\|\psi(N,x)\|. 


$$ 


Оценки $\|P_\lambda\psi\|\le c_2\|q(-i\partial/\partial x)\psi\|+ 


c\|\psi(N,x)\|$ в силу (5.15), (5.25) почти очевидны. 


 


В случае $C_\alpha(x)\in C^{\{\alpha\}}_{m,m}(b-a)$ имеем 


\begin{multline} 


P^+_{\ov\lambda}\psi=\sum_{\alpha\in\Phi}(-1)^{[\alpha]} 


((C_\alpha(x))^*\psi(x))^{(\alpha)}-\ov\lambda\psi(x)= 


\sum_{\alpha\in\Phi}(C_\alpha(x))^*(-1)^{[\alpha]} 


\psi^{(\alpha)}(x) + \\ 


% 


+\sum_{\alpha\in\Phi}\sum_{\beta<\alpha}\wt C_{\alpha,\beta}(x) 


\psi^{(\beta)}(x)\equiv P^+(x,\partial/\partial x)\psi- 


\ov\lambda\psi. 


\end{multline} 


Почти очевидно, что для полинома $P^+(x,\xi)$ выполняются аналоги 


предположений теоремы 5.2, а также аналог оценок (5.17). В этом 


случае, как доказано выше, при достаточно малых $\varepsilon_1$, 


$\varepsilon_2$ выполняются оценки (2.8). $\quad\square$ 


\medskip 


 


{\bf Доказательство следствия 5.2.} 


Пусть $v(x)\in H$, $\psi\in L_\varphi$. Тогда в силу (5.35) 


$$ 


S_6\equiv\la(P^+_{\ov\lambda}(Q^+)^{-1}-(Q^+)^{-1} 


P^+_{\ov\lambda})\psi,\,v(x)\ra= 


\sum_{\alpha\in\Phi}\sum_{\gamma<2\alpha}\la(Q^+)^{-1} 


\psi^{(\gamma)}(x),\,d_{\gamma,\alpha}(x)(Q^{-1}v(x))\ra, 


$$ 


причем функции $d_{\gamma,\alpha}(x)\in C^{\{\gamma\}}_{m,m}(b-a)$ не 


зависят от $\psi(x)$ и $v(x)$. Напомним, что $\la\,,\,\ra$ --- 


символ скалярного произведения в $L^2_m((a,b))$. Интегрируя в последних 


равенствах по частям и учитывая $(b{-}a)$-периодичность используемых 


функций, получим 


$$ 


S_6=\sum_{\alpha\in\Phi}


\sideset{}{'}\sum\la(Q^+)^{-1}\psi^{(\beta)}(x), 


\,h_{\beta,\wt\beta,\alpha}(x)(Q^{-1}v(x))^{(\wt\beta)}\ra. 


$$ 


Здесь и ниже штрих указывает на то, что суммирование производится по 


всем мультииндексам $\beta$, $\wt\beta$, удовлетворяющим условиям 


$\beta<\alpha$, $\wt\beta\le\alpha$. В силу ограниченности функций 


$h_{\beta,\wt\beta,\alpha}(x)$, не зависящих от $\psi(x)$, $v(x)$, 


имеем 


$$ 


|S_6|\le B\sum_{\alpha\in\Phi}\sideset{}{'}\sum\|(Q^+)^{-1} 


\psi^{(\beta)}(x)\|\,\|(Q^{-1}v(x))^{(\wt\beta)}\|. 


$$ 


Так как $\beta<\alpha$, $\alpha\in\Phi$ и имеет место (5.16), то 


\begin{equation} 


\|(Q^+)^{-1}\psi^{(\beta)}(x)\|\le\varepsilon_1B\|\psi(x)\|+ 


c\|\psi(N,x)\|. 


\end{equation} 


Из (5.16) и из неравенства $\wt\beta\le\alpha$ следует 


$\|(Q^{-1}v(x))^{(\wt\beta)}\|\le g\|v(x)\|$. Таким образом, доказано 


$$ 


|S_6|\le(g_3\varepsilon_1\|\psi(x)\|+ 


c\|\psi(N,x)\|)\|v(x)\|, 


$$ 


постоянные $g_3$, $\varepsilon_1$, $\varepsilon$, $N$, $c$ не зависят 


от $\psi(x)$ и от $v(x)$; постоянная $g_3$ не зависит от 


$\varepsilon_1$. В силу произвольности $v(x)$ имеет место (2.13). 


$\quad\square$ 


\medskip 


 


{\bf Доказательство следствия 5.3.} 


Если $q(i\xi)$ --- гипоэллиптический полином, то при каждом 


$\varepsilon_1>0$, некотором $N(\varepsilon_1)$ и всех $|k|\ge 


N(\varepsilon_1)$ выполняются оценки (5.16). $\quad\square$ 


 


\begin{thms} 


Пусть выполняются предположения теоремы $5.2$, оценки $(5.17)$ и 


$\varepsilon_1$ достаточно мало. Тогда 


\begin{itemize} 


\item[1.] спектр $\varphi$-задачи для уравнения $(7)$ точечен$;$ 


\item[2.] в случае существования $\varphi$-решение уравнения $(7)$ 


принадлежит $W^{\wt\Phi,2}_m(b-a)$, где 


$\wt\Phi=\mathop{\cup}\limits_{\alpha\in\Phi}\{\gamma: 


0\le\gamma\le\alpha\};$ 


\item[3.] если, кроме того, при всех $\alpha\in\Phi$ выполняются 


включения $C_\alpha(x)\in C^{\{\wt\gamma\}}_{m,m}(b-a)$, $f(x)\in 


W^{\{\wt\gamma\},2}_m(b-a)$ и мультииндекс $\wt\gamma$, не зависящий от 


$\alpha$, удовлетворяет при некотором $\beta\in\Phi$ условию 


$(\sgn\wt\gamma_1,\dotsc,\sgn\wt\gamma_n)\le\beta$, то 


$\varphi$-решение уравнения $(7)$ при условии существования 


$\varphi$-решения принадлежит $W^{\wt\Phi+\wt\gamma,2}_m(b-a)$, где 


$\wt\Phi+\wt\gamma=\mathop{\cup}\limits_{\alpha\in\wt\Phi} 


\{\alpha+\wt\gamma\}$. 


\end{itemize} 


\end{thms} 


 


\begin{pf} 


Напомним, что $\sgn 0=0$, $\sgn t=1$ $\forall\,t>0$. В силу теоремы 5.2 


и следствий 5.1, 5.2 к оператору $P_\lambda$ применимы теоремы 2.2, 2.3 


при достаточно малом $\varepsilon_1>0$. Изучая природу спектра 


оператора $P_\lambda$, мы не должны тревожиться о предположениях 


гладкости матричных функций $C_\alpha(x)$, используемых в следствиях 


5.1, 5.2. Заменив $C_\alpha(x)$ на $C_{\alpha,\delta_1}(x)\in 


C^{+\infty}_m(b-a)$, где 


$|C_\alpha(x)-C_{\alpha,\delta_1}(x)|<\delta_1$, при достаточно малом 


$\delta_1$ добьемся желаемого. Как отмечали, обсуждая предположения 


теоремы 2.2, природа спектра при этом не изменится. Мы не уточняем 


малость $\varepsilon_1$, т.\,к. это было сделано в процессе 


доказательства теоремы 5.2. Напомним только, что в тех случаях, когда 


полином $q(i\xi)$ гипоэллиптичен, $\varepsilon_1$ всегда можно выбрать 


сколь угодно малым. В силу сказанного имеют место первое и второе 


утверждение теоремы 5.3. Труднее доказать третье утверждение. 


 


Фиксируем $j$, при котором $\wt\gamma_j\ge 1$, и оператор $\wt Q$ 


$$ 


\wt Q w=\sum_k w_k\bigg(\frac{2\pi i k_j}{b_j-a_j}+\wt M\bigg) 


\exp\bigg(2\pi i k\bullet\frac{x-a}{b-a}\bigg), 


$$ 


где $\wt M$ --- достаточно большое вещественное число. Убедимся в 


выполнении аналогов предположений теоремы 2.4. Пусть $\varphi$-решение 


уравнения (7) существует. Обозначим его через $u$. Так как $u\in 


W^{\wt\Phi,2}_m(b-a)$ и $e(j)\le\beta$ при некотором $\beta\in\Phi$, то 


$u\in D_{\wt Q}$. Поэтому можно вести речь об уравнении (2.21). т.\,е. 


об уравнении 


\begin{equation} 


[\wt Q(P(x,\partial/\partial x)-\wt\lambda E)]v=\wt Q 


f(x)-\wt\lambda\wt Q u. 


\end{equation} 


Чтобы доказать его разрешимость, достаточно убедиться в выполнении 


аналогов предположений теоремы 5.2 и аналогов оценок (5.17). Для 


рассматриваемого случая имеем 


\begin{multline*} 


\wt Q(P(x,\partial/\partial x)-\wt\lambda E)\equiv 


\wt P(x,\partial/\partial x)\equiv\sum_{\alpha\in\Phi}( 


\partial/\partial x_j C_\alpha(x)) 


(\partial/\partial x)^\alpha+ \\ 


% 


+\bigg(\sum_{\alpha\in\Phi}C_\alpha(x) 


(\partial/\partial x_j+\wt M)(\partial/\partial x)^\alpha- 


\ov\lambda(\partial/\partial x_j+\wt M)\bigg). 


\end{multline*} 


Обозначим 


$$ 


\wt q(-i\xi)\equiv\bigg(\frac{2\pi}{b_j-a_j}\xi_j+\wt M\bigg) 


q(-i\xi). 


$$ 


Очевидно, что для $\wt q(k)$ выполняется аналог оценок (5.15). Пусть 


$\delta<\alpha$ и $\alpha\in\Phi$. Тогда $\delta=\wt\delta+e(j)$, 


причем либо $\wt\delta<\beta$ при некотором $\beta\in\Phi$, либо 


$\wt\delta=\beta$ при некотором $\beta\in\Phi$. В первом случае при 


$|k|\ge N(\varepsilon_1)$ 


$$ 


|(i k)^\delta|(|\wt q(k)|^{-1})=|k_j|\,|(i k)^{\wt\delta}|\bigg( 


\bigg|-\frac{2\pi i k_j}{b_j-a_j}+\wt M\bigg|^{-1}\bigg)


(|q(k)|^{-1})\le\bigg(\frac{b_j-a_j}{2\pi}\bigg) 


\varepsilon_1\equiv M_1. 


$$ 


Во втором случае имеем 


$$ 


|(i k)^\delta|\,|(\wt q(k))^{-1}|\le\frac{1}{\wt M} 


|(i k)^\beta|\,|(q(k))^{-1}|\le\frac{c}{\wt M},\quad 


\text{если} 


\quad|k|\ge N(\varepsilon_1), 


$$ 


причем $c$ не зависит от $\wt M$ и $k$. Полагая 


$\wt\varepsilon_1=\max\{M_1,\frac{c}{\wt M}\}$, окончательно убедимся в 


выполнимости для полинома $\wt q(-i\xi)$ аналога оценок (5.16). В силу 


(5.17) при $\wt\lambda=-\alpha_0t$ имеем 


\begin{multline*} 


\bigg|\wt P\bigg(x,\frac{2\pi i k}{b-a}\bigg)z\bigg|^2=\bigg(\bigg| 


\frac{2\pi i k_j}{b_j-a_j}+\wt M\bigg|^2-B_3\bigg)


\bigg|P\bigg(\bigg(x, 


\frac{2\pi i k}{b-a}\bigg)+\alpha_0t E\bigg)z\bigg|^2\ge \\ 


% 


\ge0.5\bigg(\bigg(\frac{2\pi k_j}{b_j-a_j}\bigg)^2+\wt M^{\,2}\bigg) 


\{d_1(t)+d_2|q(k)|^2\}|z|^2 \ge \\ 


% 


\ge0.5\bigg\{d_1(t)+d_2\bigg(\bigg|\bigg(\frac{2\pi i k_j}{b_j-a_j}+ 


\wt M\bigg)q(k)\bigg|^2\bigg)\bigg\}|z|^2= 


0.5(d_1(t)+d_2|\wt q(k)|^2)|z|^2. 


\end{multline*} 


Поэтому выполняется аналог оценок (5.17). Таким образом, доказано, что 


к оператору $\wt P_0$, порожденному полиномом $\wt P(x,\xi)$, применимы 


аналоги утверждений теоремы 5.2 и следствий 5.1, 5.2, 5.3. 


Следовательно, уравнение (5.37) разрешимо. Пусть $v$ --- его 


$\varphi$-решение, т.\,е. 


$$ 


\sum_k\wt Q(P(x,\partial/\partial x)-\wt\lambda E) 


\bigg(v_k\exp\bigg(2\pi i k\bullet\frac{x-a}{b-a}\bigg)\bigg)= 


\wt Q f(x)-\wt\lambda\wt Q u, 


$$ 


и ряд сходится в $H$ по норме. Так как оператор $\wt Q$ имеет 


непрерывный обратный, то 


$$ 


\sum_k(P(x,\partial/\partial x)-\wt\lambda E) 


\bigg(v_k\exp\bigg(2\pi i k\bullet\frac{x-a}{b-a}\bigg)\bigg)= 


f(x)-\wt\lambda u, 


$$ 


причем ряд снова сходится в $H$ по норме. Отсюда следует, что $v$ --- 


$\varphi$-решение уравнения 


$(P(x,\partial/\partial x)-\wt\lambda E)v=f(x)-\wt\lambda u$, причем 


$u$ --- также его $\varphi$-решение. Однако $\wt\lambda=-\alpha_0t$, 


поэтому можно считать $\wt\lambda$ несобственным значением оператора 


$P_{\wt\lambda}$ и, следовательно, $u=v$. Так как $v\in D_{(Q\wt Q)}$, 


т.\,е. $v\in W^{\wt\Phi+e(j),2}_m(b-a)$, то $u\in 


W^{\wt\Phi+e(j),2}_m(b-a)$. Таким образом, доказана справедливость 


утверждения 3 в том случае, когда $\wt\gamma=e(j)$. Повторив 


приведенные выше рассуждения многократно, окончательно докажем теорему 


5.3. 


\end{pf} 


 


3. {\it Функции $\varphi_{(p)}$, $p\in N$, являются сужениями на 


$\Omega$ некоторых функций из $W^{\Phi,2}_m(b-a)$.} В данном пункте 


предполагается, что $H=L^2_m(\Omega)$, $\varphi=\{e(j)y_{(p)},\ 


j=1,\dotsc,m,\ p\in N\}$, и функции $y_{(p)}\equiv y_{(p)}(x)$ являются 


сужениями на $\Omega$ некоторых функций $\wt y_{(p)}\in 


W^{\Phi,2}_m(b-a)$. Выделим случаи, когда $\varphi$-решения уравнения 


(7), в случае их существования, также будут сужениями на $\Omega$ 


некоторых $(b{-}a)$-периодических решений соответствующих уравнений. 


 


Отмеченная ситуация типична для многих линейных задач, особенно в 


случае $\Omega=(\wt a,\wt b)\subset[a,b]$ (см. cс.\,30, 31 из 


\cite{1}). Предположим, что открытое множество $\Omega\subset[a,b]$ имеет 


такую границу, что имеет смысл запись $C^{+\infty}_0(\Omega)$, причем 


при некоторых $h_{(k,r)}(x)\in C^{+\infty}_0(\Omega)$ и каждом $k\in N$ 


\begin{equation} 


\sum_{\alpha\in\Phi\cup\{0\}}\|h^{(\alpha)}_{(k,r)}(x)- 


y^{(\alpha)}_{(k)}(x)\|\to 0\quad 


\text{при}\quad r\to+\infty. 


\end{equation} 


Тогда при каждом $k$ функции $y_{(k)}$ являются сужениями на $\Omega$ 


некоторых функций $\wt y_{(k)}{\in}W^{\Phi,2}_1(b{-}a)$. Убедимся в этом. 


Очевидно, каждая функция $h_{(k,r)}(x)$ является сужением на $\Omega$ 


некоторой функции 


$\wt h_{(k,r)}(x)\in W^{+\infty,2}_1(b-a)$. В силу (5.38) имеем 


$$ 


\sum_{\alpha\in\Phi\cup\{0\}}\|\wt h^{(\alpha)}_{(k,r)}(x)- 


\wt h^{(\alpha)}_{(k,v)}(x)\|_1\to 0\quad 


\text{при}\quad r\ge v\to+\infty, 


$$ 


где $\|\cdot\|_1$ --- символ нормы в $L^2_1(b-a)$. Отсюда следует, что 


для каждого $k$ последовательность


$\{\wt h_{(k,\nu)}(x)\}$ сходится к 


некоторой функции $\wt y_{(k)}(x)\in W^{\Phi,2}_1(b-a)$. Очевидно, на 


$\Omega$ функции $y_{(k)}(x)$, $\wt y_{(k)}(x)$ совпадают. 


 


Наряду с $\varphi$ рассмотрим последовательность $\{e(j)\wt 


y_{(k)}(x),\ j=1,\dotsc,m,\ k\in N\}=\wt\varphi$, полином $\wt 


P(x,z)=\sum\limits_{\alpha\in\Phi}\wt C_\alpha(x)(z^\alpha)$, в котором 


$\wt C_\alpha(x)\in L^2_{m,m}(b-a)$ совпадают на $\Omega$ с 


$C_\alpha(x)$, причем элементы матриц $\wt C_\alpha(x)$ ограничены в 


существенном в $[a,b]$, 


 


\begin{thms} 


Предположим, что $\wt\varphi$-задача для уравнения 


\begin{equation} 


(\wt P(x,\partial/\partial x)-\lambda E)\wt v=\wt f(x), 


\ \; x\in(a,b);\quad \wt f(x)=f(x)\ \;\text{при}\ \;x\in\Omega 


\end{equation} 


имеет при некоторой $\wt f(x)\in L^2_m(b-a)$ решение $\wt v$. Тогда 


сужение на $\Omega$\; $\wt\varphi$-решения $\wt v$ является 


$\varphi$-решением уравнения $(7)$. 


\end{thms} 


 


\begin{pf} 


Последовательность $\wt\varphi$ имеет биортогональную, ибо 


биортогональную имеет $\varphi$ и $\Omega\subset[a,b]$. Так как $\wt v$ 


--- $\wt\varphi$-решение уравнения (5.39), то 


$$ 


\sum_{k,j}v_{k,j}(\wt P(x,\partial/\partial x)-\lambda E) 


(e(j)\wt y_{(k)}(x))=\wt f(x), 


$$ 


причем ряд сходится в $H$ по норме. Рассмотрим $\varphi$-распределение 


$v$, определяемое коэффициентами Фурье $v_{k,j}$, $j=1,\dotsc,m$, $k\in 


N$ (см. замечания 1.2, 1.3). Именно это $\varphi$-распределение будем 


называть сужением на $\Omega$\; $\wt\varphi$-распределения $\wt v$. 


Легко убедиться в том, что в регулярном случае, т.\,е. в случае $\wt 


v\in L^2_m(b-a)$, предложенное понятие совпадает с классическим 


понятием сужения функций. Так как $\Omega\subset[a,b]$, то последний 


ряд сходится и в $L^2_m(\Omega)$ по норме, причем его сумма в 


$L^2_m(\Omega)$ совпадает с $f(x)$. А это означает, что $v$ --- 


$\varphi$-решение уравнения (7). 


\end{pf} 


 


\begin{notes} 


Если известно, что $\varphi$-задача для уравнения (7) не имеет двух 


решений, и $\wt\varphi$-задача для уравнения (5.39) разрешима, что 


$\varphi$-решение уравнения (7) существует, единственно и является 


сужением на $\Omega$\; $\wt\varphi$-решений уравнения (5.39). 


\end{notes} 


Обозначим через $\wt P_\lambda$ оператор, порожденный 


$\wt\varphi$-задачей для (5.39); через $\wt{\wt P}_\lambda$ --- 


оператор, порожденный $\wt{\wt\varphi}$-задачей для (5.39). Пусть 


$\wt{\wt\varphi}$ совпадает с (1.4). Другими словами, 


$\wt{\wt\varphi}$-задача для уравнения (5.39) является 


$(b{-}a)$-периодической задачей, исследованной в п.\,2. 


 


\begin{thms} 


Предположим, что $D_{\wt P_\lambda}=D_Q$, $D_{\wt{\wt 


P}_\lambda}=D_{\wt Q}$ и операторы $Q$, $\wt Q$ соответственно 


$\wt\varphi$-, $\wt{\wt\varphi}$-стационарные и взаимно однозначные. Тогда 


в случае существования $\wt\varphi$-решение $\wt v$ уравнения $(5.39)$ 


продолжается до $(b{-}a)$-периодического решения $\wt{\wt v}$ того же 


уравнения$;$ если, кроме того, $\wt\varphi$ --- базис в $L^2_m((a-b))$ 


и $D_{\wt Q}\subset L^2_m(b-a)$, то $\wt v=\wt{\wt v}$ как элементы 


пространства $L^2_m(b-a)$. 


\end{thms} 


 


\begin{pf} 


Так как $\wt\varphi_{(p)}(x)\in W^{\Phi,2}_m(b-a)$, то 


$$ 


\wt\varphi^{\,(\alpha)}_{(p)}(x)=\sum_{q_0} 


a_{p,q_0}\bigg(\frac{2\pi i q_0}{b-a}\bigg)^\alpha\exp\bigg( 


2\pi i q_0\bullet\frac{x-a}{b-a}\bigg)\ \; \forall\,\alpha\in\Phi, 


$$ 


причем ряды сходятся в $L^2_m((a,b))$ по норме. С другой стороны, 


элементы матриц $\wt C_\alpha(x)$ ограничены в существенном, поэтому 


$$ 


\wt P_\lambda\wt\varphi_{(p)}(x)=\sum_{q_0}(\wt P 


(x,\partial/\partial x)-\lambda E)\bigg(a_{p,q_0}\exp\bigg( 


2\pi i q_0\bullet\frac{x-a}{b-a}\bigg)\bigg), 


$$ 


и ряд сходится в $L^2_m((a,b))$ по норме. Последние равенства вместе с 


остальными предположениями теоремы гарантируют выполнимость аналогов 


предположений теоремы 3.1. Из последней следует, что $\wt v$ 


продолжается до $(b{-}a)$-периодического решения $\wt{\wt v}$ уравнения 


(5.39). Однако $\wt{\wt v}\in L^2_m((a,b))$, поэтому $\wt v\in 


L^2_m((a,b))$ и $(\wt v-\wt{\wt v})(\wt\varphi_{(p)}(x))=0$. В 


частности, если $\wt\varphi$ --- базис в $L^2_m((a,b))$, то $\wt 


v=\wt{\wt v}$ как элементы пространства $L^2_m((a,b))$. 


\end{pf} 


 


\begin{cors} 


Пусть выполняются предположения теоремы 5.5, последовательность 


$\wt\varphi$ является базисом в $L^2_m((a,b))$, и уравнение (5.39) 


имеет $\wt\varphi$-решение. Тогда одно из $\varphi$-решений уравнения 


(7) является сужением на $\Omega$\; $(b{-}a)$-периодического решения 


уравнения (5.39). Если, кроме того, уравнение (7) не может иметь двух 


$\varphi$-решений, то $\varphi$-решение уравнения (7) является сужением 


на $\Omega$\; $(b{-}a)$-периодического решения уравнения (5.39). 


\end{cors} 


 


\begin{pf} 


Пусть $\wt v$ --- $\wt\varphi$-решение уравнения (5.39). В силу теоремы 


5.5 оно продолжается до $(b{-}a)$-периодического решения $\wt{\wt v}$ 


уравнения (5.39). Итак, $\wt v=\wt{\wt v}$ как $\varphi$-распределения. 


Однако сужение на $\Omega$\; $\wt\varphi$-распределения $\wt v$ (а 


значит, и $\wt{\wt v}$), как доказано в теореме 5.4, является 


$\varphi$-решением уравнения (7). Первое утверждение доказано. В силу 


первого утверждения второе почти очевидно. 


\end{pf} 


 


\begin{notes} 


Понятно, что из результатов части III следует теория периодической 


задачи, изложенная в \cite{4}. Обратное невозможно в принципе. 


Действительно, если уравнение $P(x,\partial/\partial x)u=f(x)$ имеет 


$(b{-}a)$-периодическое решение, и $\wt{\wt P}(\partial/\partial x)$ 


--- линейный дифференциальный квазиполином с постоянными 


коэффициентами, удовлетворяющий условию \linebreak


$\det(\wt{\wt P}(2\pi i\frac{k}{b-a}))\ne0$ $\forall\,k\in\Bbb Z^n$, то 


и уравнение $P(x,\partial/\partial x)\wt{\wt P}(\partial/\partial 


x)v=f(x)$ имеет $(b{-}a)$-перио\-ди\-чес\-кое решение. Понятие периодического 


решения, использованное в \cite{4}, делает отмеченное утверждение 


невозможным. В этом мы убедились, рассматривая математическую модель 


колебаний струны. Аналогичное можно сказать и о других 


$\varphi$-задачах, в том числе и о задаче, рассмотренной в ([2], п.\,1).


Для последней вместо $\wt{\wt P}(\partial/\partial x)$ следует 


взять $\wt{\wt P}((\partial/\partial x)^2+h)$ и предположить 


$\det(\wt{\wt P}(\Lambda_k))\ne0$ $\forall\,k\in\Bbb Z^n$. В тех 


случаях, когда $\Phi=\{\alpha:|\alpha|\le n_1\}$, $m=1$, и (5.17) имеет 


место при всех $k\in\Bbb R^n$, в \cite{5} дифференциальный полином 


$P(x,\partial/\partial x)$ называется оператором главного типа. 


\end{notes} 


 


 


\begin{thebibliography}{99} 


 


\bibitem{1} 


Мокейчев В.С., \fbox{Мокейчев А.В.} {\em Новый подход к теории линейных 


задач для систем дифференциальных уравнений в частных производных}. I 


// Изв. вузов. Математика. -- 1999. -- \No\,1. -- С.\,25--35. 


\bibitem{2} 


Мокейчев В.С., \fbox{Мокейчев А.В.} {\em Новый подход к теории линейных 


задач для систем дифференциальных уравнений в частных производных}. II


// Изв. вузов. Математика. -- 1999. -- \No\,7. -- С.\,30--41. 


\bibitem{3} 


Мокейчев В.С. {\em Дифференциальные уравнения с отклоняющимися 


аргументами}. -- Казань: Изд-во Казанск. ун-та, 1985. -- 222 с. 


\bibitem{4} 


Берс Л., Джон Ф., Шехтер М. {\em Уравнения с частными производными}. -- 


М.: Мир, 1966. -- 357~c. 


\bibitem{5} 


Егоров Ю.В. {\em Линейные дифференциальные уравнения главного типа}. -- 


М.: Наука, 1984. -- 359 с. 


 


\end{thebibliography} 


 


\vskip 1cm 


 


\post{\it Казанский государственный}{\it Поступила} 


\post{\it университет}{18.03.1996} 


 


\label{end} 


\end{document} 


 


